\color{unidarkgreen}{\subsection{Дифференциальная токовая защита}}
\color{black}

\begin{enumerate}

\item\textbf{Общие сведения}

\begin{enumerate}

\item 
Дифференциальная токовая защита (ДТЗ) применяется в качестве основной быстродействующей защиты силового оборудования (трансформаторы, реакторы и др.) от всех видов КЗ в обмотках и на выводах. Цепи измерения защиты подключаются на фазные токи контролируемых сторон. Защита может контролировать до шести трехфазных цепей измерения. За положительное направление токов для всех контролируемых сторон принято направление к защищаемому объекту.

\item 
Параметры для настройки ДТЗ приведены в таблице \ref{app_set:dzt}.

\end{enumerate}

\item\textbf{Расчетные величины}

\begin{enumerate}

\item 
Перед расчетом дифференциальных токов и тока торможения производится цифровое выравнивание токов по амплитуде и компенсация фазового сдвига.

\item
Цифровое выравнивание по амплитуде обеспечивается умножением измеренного значения тока на соответствующий коэффициент, определяемый по формуле: 

\begin{equation}
\label{eq:Kampl}
K_{\textsl{ампл, N}} = \frac{I_{\textsl{ном.ТТ, перв,N}}}{I_{\textsl{баз,N}} \cdot K_{\textsl{сх}}},
\end{equation}

\begin{eqexpl}[10mm]
\eqexplSetItemAlign{r}  % выравнивание обозначений по левому краю(l)/правому краю(r)/центру(c)
\item{$K_{\textsl{сх}}$} коэффициент схемы, равный 1,0 при сборке измерительных ТТ в звезду и $\sqrt{3}$ при сборке ТТ в треугольник;
\item{$I_{\textsl{ном.ТТ, перв.N}}$} номинальный первичный ток измерительного ТТ стороны N;
\item{$I_{\textsl{баз,N}}$} базисный ток стороны N, значение которого рассчитывается по формуле:
\end{eqexpl}

\begin{equation}
\label{eq:Ibaz}
I_{\textsl{баз,N}} = \frac{S_{\textsl{баз}}}{\sqrt{3}U_{\textsl{ном,N}}},
\end{equation}

\begin{eqexpl}[10mm]
\eqexplSetItemAlign{r}  % выравнивание обозначений по левому краю(l)/правому краю(r)/центру(c)
\item{$U_{\textsl{ном.N}}$} номинальное напряжение стороны N силового трансформатора;
\item{$S_{\textsl{баз}}$} базисная мощность.
\end{eqexpl}

\item 
За базисную мощность для всех сторон принимается мощность защищаемого первичного оборудования. Если обмотки первичного оборудования имеют разную мощность, то за базисную мощность для всех сторон принимается мощность наиболее мощной обмотки.

\item
Компенсация фазового сдвига для каждой из сторон выполняется индивидуально по формулам, приведенным в таблице \ref{dtz:tbl1}. Для обеспечения нечувствительности дифференциальной защиты к внешним однофазным коротким замыканиям, когда токи нулевой последовательности могут протекать через защиту и при этом восприниматься как дифференциальный ток, производится устранение токов нулевой последовательности. Устранение тока нулевой последовательности необходимо в случаях, когда ток нулевой последовательности может протекать только по одну сторону защищаемого объекта. 

\item
После предварительной цифровой обработки измерений производится расчет дифференциальных и тормозного токов. Дифференциальные токи фаз определяются по формуле:

\begin{equation}
\label{eq:Idiff1}
i_{\textsl{диф.x}} = \sum_{N=1}^k i'_{\textsl{x,N}} \cdot K_{\textsl{ампл,N}}
\end{equation}

\begin{eqexpl}[5mm]
\eqexplSetItemAlign{r}  % выравнивание обозначений по левому краю(l)/правому краю(r)/центру(c)
\item{x} наименование фазы (A, B, C);
\item{k} количество контролируемых плеч
\end{eqexpl}

\item
Ток торможения является общим для всех фаз и численно равен току, имеющему максимальное действующее значение. 

\begin{equation}
\label{eq:Idiff2}
I_{\textsl{торм}} = \max_{N=1..k} \left[RMS(i'_{\textsl{x,N}}) \cdot K_{\textsl{ампл,N}}\right]
\end{equation}


%%%%%%%%%%%%%%%%%%%%%%%%%%%%%%%%%%%%%%%%%%%%%%%%%%%%%%%%%%%%%%%%%%%%%%%%%% ТАБЛИЦА 1 %%%%%%%%%%%%%%%%%%%%%%%%%%%%%%%%%%%%%%%%%%%%%%%%%%%%%%%%%%%%%%%%%%%
\newcolumntype{A}{>{\centering\arraybackslash}m{1cm}}
\newcolumntype{B}{>{\centering\arraybackslash}m{4.07cm}}
\setlength{\extrarowheight}{3pt}  % задание высоты строк
\footnotesize

\needspace{10\baselineskip} % ЗДЕСЬ ВАЖНО

\begin{longtable}{|A|B|A|B|A|B|} % здесь в скобках указываются количество столбцов, их параметры и вертикальные линии
\caption{Формулы компенсации фазового сдвига \hfill\vspace{-0.5\baselineskip}}\label{dtz:tbl1}\\


\hline
\rowcolor{gray!30} № Сх. & Формула & № Сх. & Формула & № Сх. & Формула \\
\hline % горизонтальная линия (вся строка)
\endfirsthead

\caption*{\hspace{3pt}\emph{Продолжение таблицы \ref{dtz:tbl1}\hfill\vspace{-0.5\baselineskip}}} \\ % сделано по ГОСТ 2.105 п.6.8.7

\hline % горизонтальная линия (вся строка)
\rowcolor{gray!30} № Сх. & Формула & № Сх. & Формула & № Сх. & Формула \\
\hline % горизонтальная линия (вся строка)
\endhead
\hline % горизонтальная линия (вся строка)
\endfoot
\endlastfoot

\multirow{3}{*}{0} & $i'_{A} = i_{A}$ & \multirow{3}{*}{8} & $i'_{A} = i_{C}$ & \multirow{3}{*}{16} & $i'_{A} = i_{B}-(i_{A}+i_{B}+i_{C})/3$ \\
\cline{2-2} \cline{4-4} \cline{6-6} % горизонтальная линия только для столбцов 2, 4, 6
& $i'_{B} = i_{B}$ &  & $i'_{B} = i_{A}$ & & $i'_{B} = i_{C}-(i_{A}+i_{B}+i_{C})/3$  \\
\cline{2-2} \cline{4-4} \cline{6-6} % горизонтальная линия только для столбцов 2, 4, 6
& $i'_{C} = i_{C}$ &  & $i'_{C} = i_{B}$ & & $i'_{C} = i_{A}-(i_{A}+i_{B}+i_{C})/3$  \\
\hline % горизонтальная линия (вся строка)
\multirow{3}{*}{1} & $i'_{A} = (i_{A}-i_{C})/\sqrt3$ & \multirow{3}{*}{9} & $i'_{A} = (i_{C}-i_{B})/\sqrt3$ & \multirow{3}{*}{17} & $i'_{A} = (i_{B}-i_{A})/\sqrt3$ \\
\cline{2-2} \cline{4-4} \cline{6-6} % горизонтальная линия только для столбцов 2, 4, 6
& $i'_{B} = (i_{B}-i_{A})/\sqrt3$ &  & $i'_{B} = (i_{A}-i_{C})/\sqrt3$ & & $i'_{B} = (i_{C}-i_{B})/\sqrt3$  \\
\cline{2-2} \cline{4-4} \cline{6-6} % горизонтальная линия только для столбцов 2, 4, 6
& $i'_{C} = (i_{C}-i_{B})/\sqrt3$ &  & $i'_{C} = (i_{B}-i_{A})/\sqrt3$ & & $i'_{C} = (i_{A}-i_{C})/\sqrt3$  \\
\hline % горизонтальная линия (вся строка)
\multirow{3}{*}{2} & $i'_{A} = -i_{C}$ & \multirow{3}{*}{10} & $i'_{A} = -i_{B}$ & \multirow{3}{*}{18} & $i'_{A} = -i_{A}+(i_{A}+i_{B}+i_{C})/3$ \\
\cline{2-2} \cline{4-4} \cline{6-6} % горизонтальная линия только для столбцов 2, 4, 6
& $i'_{B} = -i_{A}$ &  & $i'_{B} = -i_{C}$ & & $i'_{B} = -i_{B}+(i_{A}+i_{B}+i_{C})/3$  \\
\cline{2-2} \cline{4-4} \cline{6-6} % горизонтальная линия только для столбцов 2, 4, 6
& $i'_{C} = -i_{B}$ &  & $i'_{C} = -i_{A}$ & & $i'_{C} = -i_{C}+(i_{A}+i_{B}+i_{C})/3$  \\
\hline % горизонтальная линия (вся строка)
\multirow{3}{*}{3} & $i'_{A} = (i_{B}-i_{C})/\sqrt3$ & \multirow{3}{*}{11} & $i'_{A} = (i_{A}-i_{B})/\sqrt3$ & \multirow{3}{*}{19} & $i'_{A} = (i_{C}-i_{A})/\sqrt3$ \\
\cline{2-2} \cline{4-4} \cline{6-6} % горизонтальная линия только для столбцов 2, 4, 6
& $i'_{B} = (i_{C}-i_{A})/\sqrt3$ &  & $i'_{B} = (i_{B}-i_{C})/\sqrt3$ & & $i'_{B} = (i_{A}-i_{B})/\sqrt3$  \\
\cline{2-2} \cline{4-4} \cline{6-6} % горизонтальная линия только для столбцов 2, 4, 6
& $i'_{C} = (i_{A}-i_{B})/\sqrt3$ &  & $i'_{C} = (i_{C}-i_{A})/\sqrt3$ & & $i'_{C} = (i_{B}-i_{C})/\sqrt3$  \\
\hline % горизонтальная линия (вся строка)
\multirow{3}{*}{4} & $i'_{A} = i_{B}$ & \multirow{3}{*}{12} & $i'_{A} = i_{A}-(i_{A}+i_{B}+i_{C})/3$ & \multirow{3}{*}{20} & $i'_{A} = i_{C}-(i_{A}+i_{B}+i_{C})/3$ \\
\cline{2-2} \cline{4-4} \cline{6-6} % горизонтальная линия только для столбцов 2, 4, 6
& $i'_{B} = i_{C}$ &  & $i'_{B} = i_{B}-(i_{A}+i_{B}+i_{C})/3$ & & $i'_{B} = i_{A}-(i_{A}+i_{B}+i_{C})/3$  \\
\cline{2-2} \cline{4-4} \cline{6-6} % горизонтальная линия только для столбцов 2, 4, 6
& $i'_{C} = i_{A}$ &  & $i'_{C} = i_{C}-(i_{A}+i_{B}+i_{C})/3$ & & $i'_{C} = i_{B}-(i_{A}+i_{B}+i_{C})/3$  \\
\hline % горизонтальная линия (вся строка)
\multirow{3}{*}{5} & $i'_{A} = (i_{B}-i_{A})/\sqrt3$ & \multirow{3}{*}{13} & $i'_{A} = (i_{A}-i_{C})/\sqrt3$ & \multirow{3}{*}{21} & $i'_{A} = (i_{C}-i_{B})/\sqrt3$ \\
\cline{2-2} \cline{4-4} \cline{6-6} % горизонтальная линия только для столбцов 2, 4, 6
& $i'_{B} = (i_{C}-i_{B})/\sqrt3$ &  & $i'_{B} = (i_{B}-i_{A})/\sqrt3$ & & $i'_{B} = (i_{A}-i_{C})/\sqrt3$  \\
\cline{2-2} \cline{4-4} \cline{6-6} % горизонтальная линия только для столбцов 2, 4, 6
& $i'_{C} = (i_{A}-i_{C})/\sqrt3$ &  & $i'_{C} = (i_{C}-i_{B})/\sqrt3$ & & $i'_{C} = (i_{B}-i_{A})/\sqrt3$  \\
\hline % горизонтальная линия (вся строка)
\multirow{3}{*}{6} & $i'_{A} = -i_{A}$ & \multirow{3}{*}{14} & $i'_{A} = -i_{C}+(i_{A}+i_{B}+i_{C})/3$ & \multirow{3}{*}{22} & $i'_{A} = -i_{B}+(i_{A}+i_{B}+i_{C})/3$ \\
\cline{2-2} \cline{4-4} \cline{6-6} % горизонтальная линия только для столбцов 2, 4, 6
& $i'_{B} = -i_{B}$ &  & $i'_{B} = -i_{A}+(i_{A}+i_{B}+i_{C})/3$ & & $i'_{B} = -i_{C}+(i_{A}+i_{B}+i_{C})/3$  \\
\cline{2-2} \cline{4-4} \cline{6-6} % горизонтальная линия только для столбцов 2, 4, 6
& $i'_{C} = -i_{C}$ &  & $i'_{C} = -i_{B}+(i_{A}+i_{B}+i_{C})/3$ & & $i'_{C} = -i_{A}+(i_{A}+i_{B}+i_{C})/3$  \\
\hline % горизонтальная линия (вся строка)
\multirow{3}{*}{7} & $i'_{A} = (i_{C}-i_{A})/\sqrt3$ & \multirow{3}{*}{15} & $i'_{A} = (i_{B}-i_{C})/\sqrt3$ & \multirow{3}{*}{23} & $i'_{A} = (i_{A}-i_{B})/\sqrt3$ \\
\cline{2-2} \cline{4-4} \cline{6-6} % горизонтальная линия только для столбцов 2, 4, 6
& $i'_{B} = (i_{A}-i_{B})/\sqrt3$ &  & $i'_{B} = (i_{C}-i_{A})/\sqrt3$ & & $i'_{B} = (i_{B}-i_{C})/\sqrt3$  \\
\cline{2-2} \cline{4-4} \cline{6-6} % горизонтальная линия только для столбцов 2, 4, 6
& $i'_{C} = (i_{B}-i_{C})/\sqrt3$ &  & $i'_{C} = (i_{A}-i_{B})/\sqrt3$ & & $i'_{C} = (i_{C}-i_{A})/\sqrt3$  \\
\hline % горизонтальная линия (вся строка)
\end{longtable}
\normalsize

\end{enumerate}

\item\textbf{Измерительные органы}

\begin{enumerate}

\item
Функция содержит две ступени дифференциальной защиты: чувствительную (с торможением) и грубую (отсечку). Ступень отсечки предназначена для быстрого отключения защищаемого объекта без учета критериев торможения при тяжелых внутренних повреждениях, сопровождающихся большим дифференциальным током. Срабатывание данной ступени осуществляется при превышении дифференциальным током любой из фаз значения \textbf{«Iср ДТО»}. Алгоритм работы ДТО приведен на рисунке \ref{pict:DTO}.

\vspace{2mm}
\begin{figure}[H]
\centering
\includegraphics[width=1\textwidth,height=1\textheight,keepaspectratio]{\fbpath/06.02 ДЗАТ/041.0611 ДТЗ/img/DTO.pdf}
\captionof{figure}{Функциональная схема логики дифференциальной токовой отсечки}\label{pict:DTO}
\end{figure}

\item
Ступень дифференциальной защиты с торможением предназначена для защиты от повреждений, сопровождающихся малыми дифференциальными токами, и имеет характеристику срабатывания, состоящую из трех участков (рис.\ref{pict:dz_har}).

\begin{figure}[H]
\centering
\includegraphics[width=0.5\textwidth,height=0.5\textheight,keepaspectratio]{\fbpath/06.02 ДЗАТ/041.0611 ДТЗ/img/dz_har.pdf}
\captionof{figure}{Характеристика срабатывания дифференциальной защиты}\label{pict:dz_har}
\end{figure}

\item
Характеристика торможения дифференциальной защиты задается следующими параметрами:
\begin{itemize}
    \item \textbf{Iср диф. } – минимальный дифференциальный ток срабатывания;
    \item \textbf{Iт1} – начальный ток торможения второго участка;
    \item \textbf{Кт2} – коэффициент наклона второго участка;
    \item \textbf{Iт2} – начальный ток торможения третьего участка;
    \item \textbf{Кт2} – коэффициент наклона третьего участка.
\end{itemize}

\item
Выходной сигнал срабатывания характеристики торможения на соответствующем участке определяется по формулам: 

\begin{itemize}
    \item участок 1: $I_{\textsl{Д}} \geq I_{\textsl{ср.диф}}$, при $I_{\textsl{Т}} \leq I_{\textsl{Т1}}$
    \item участок 2: $I_{\textsl{Д}} \geq I_{\textsl{ср.диф}} + K_{\textsl{Т1}} \cdot (I_{\textsl{Т}} - I_{\textsl{Т1}})$, при $I_{\textsl{Т1}} < I_{\textsl{Т}} \leq I_{\textsl{Т2}}$
    \item участок 3: $I_{\textsl{Д}} \geq I_{\textsl{ср.диф}} + K_{\textsl{Т1}} \cdot (I_{\textsl{Т2}} - I_{\textsl{Т1}}) + K_{\textsl{Т2}} \cdot (I_{\textsl{Т}} - I_{\textsl{Т2}}) $, при $I_{\textsl{Т}} > I_{\textsl{Т2}}$
\end{itemize}

\item
Алгоритм работы ДТО приведен на рисунке \ref{pict:DTZ}.

\vspace{2mm}
\begin{figure}[H]
\centering
\includegraphics[width=1\textwidth,height=1\textheight,keepaspectratio]{\fbpath/06.02 ДЗАТ/041.0611 ДТЗ/img/DTZ.pdf}
\captionof{figure}{Функциональная схема логики дифференциального токового органа с торможением}\label{pict:DTZ}
\end{figure}

\end{enumerate}

\item\textbf{Блокировка по второй гармонике}

\begin{enumerate}

\item
Причиной возникновений броска тока намагничивания в силовом оборудовании (трансформаторах, шунтирующих реакторах) является резкое изменение уровня напряжения намагничивания, характерное для режимов постановки оборудования под напряжение, восстановления уровня напряжения после отключения внешнего КЗ и др. Обнаружение режима броска тока намагничивания основано на контроле отношения величины второй гармоники дифференциального тока, как наиболее характерной для режима броска тока намагничивания, к величине составляющей первой гармоники. 

\item
Ввод блокировки в работу осуществляется программной накладкой \textbf{«Блок.2г.ДТЗт»}. Уровень второй гармоники, при котором производится блокировка дифференциальной защиты с торможением, задается значением \textbf{«I2г/I1г ДТЗт»}. 

\item
Блокировка по второй гармонике выполнена пофазной. В некоторых случаях при броске тока намагничивания содержание второй гармоники в одной из фаз может быть недостаточно большим, что может вызвать ложное срабатывание защиты. Для данных случаев может быть использована перекрестная блокировка по 2ой гармонике. Перекрестная блокировка по 2ой гармонике действует на блокирование всех фаз при превышении заданного уровня второй гармоники в одной или в двух фазах в зависимости от положения программной накладки \textbf{«Перекр.блок.2г»}. Перекрестная блокировка по второй гармонике выполняется в течение ограниченного времени, определяемого параметром \textbf{«Тперекр.блок.2г»}. Алгоритм работы Д2Г приведен на рисунке \ref{pict:d2g}.

\vspace{2mm}
\begin{figure}[H]
\centering
\includegraphics[width=1\textwidth,height=1\textheight,keepaspectratio]{\fbpath/06.02 ДЗАТ/041.0611 ДТЗ/img/d2g.pdf}
\captionof{figure}{Функциональная схема логики блокировки по второй гармонике}\label{pict:d2g}
\end{figure}

\end{enumerate}

\item\textbf{Блокировка по пятой гармонике}

\begin{enumerate}

\item
Причиной возникновений повышенной индукции силовых трансформаторов является повышение напряжения и/или понижение уровня частоты. Повышение индукции выше номинального значения быстро ведет к насыщению стального сердечника и большим потерям от вихревых токов. Обнаружение режима перенасыщения основано на контроле отношения величины составляющей пятой гармоники дифференциального тока, как наиболее характерной для режима перенасыщения, к величине составляющей первой гармоники. 

\item
Ввод блокировки в работу осуществляется программной накладкой \textbf{«Блок.5г.ДТЗт»}. Уровень пятой гармоники, при котором производится блокировка дифференциальной защиты с торможением, задается значением \textbf{«I5г/I1г ДТЗт»}. 

\item
Блокировка по пятой гармонике выполнена пофазной. В некоторых случаях содержание пятой гармоники в одной из фаз может быть недостаточно большим, что может вызвать ложное срабатывание защиты. Для данных случаев может быть использована перекрестная блокировка по 5ой гармонике. Перекрестная блокировка по 5ой гармонике действует на блокирование всех фаз при превышении заданного уровня пятой гармоники в одной или в двух фазах в зависимости от положения программной накладки \textbf{«Перекр.блок.5г»}. Перекрестная блокировка по 5ой гармонике выполняется в течение ограниченного времени, определяемого параметром \textbf{«Тперекр.блок.5г»}. Алгоритм работы Д2Г приведен на рисунке \ref{pict:d5g}.

\vspace{2mm}
\begin{figure}[H]
\centering
\includegraphics[width=1\textwidth,height=1\textheight,keepaspectratio]{\fbpath/06.02 ДЗАТ/041.0611 ДТЗ/img/d5g.pdf}
\captionof{figure}{Функциональная схема логики блокировки по пятой гармонике}\label{pict:d5g}
\end{figure}

\end{enumerate}

\item\textbf{Контроль цепей тока}

\begin{enumerate}

\item
Функция содержит логику контроля исправности измерительных цепей тока. Обнаружение неисправности измерительных цепей осуществляется при превышении дифференциальным током в любой из фаз значения \textbf{«Iнеисп ДТЗ»}. По истечении выдержки времени \textbf{«Тнеисп ДТЗ»} формируется сигнал о неисправности цепей измерения защиты. Алгоритм работы КЦТ приведен на рисунке \ref{pict:KCT}.

\vspace{2mm}
\begin{figure}[H]
\centering
\includegraphics[width=1\textwidth,height=1\textheight,keepaspectratio]{\fbpath/06.02 ДЗАТ/041.0611 ДТЗ/img/KCT.pdf}
\captionof{figure}{Функциональная схема логики контроля цепей тока}\label{pict:KCT}
\end{figure}


\end{enumerate}

\end{enumerate}